% GENERATED FILE. Edit canonical lessons, then run npm run generate:books.
% canonical-source-hash: fnv1a64:41e7c34f9be0de70
% canonical-lessons: GE-C105-restaurant, GE-C105-baeckerei, GE-C105-kino, GE-C105-theater, GE-C105-rathaus

\chapter{Restaurant, Bakery, Cinema, Theatre, Town hall}
\label{ch:ge-restaurant-bakery-cinema-theatre-town-hall}
\hlchaptermodality{105}

\begin{chapteropening}
\textbf{By the end of this chapter:} \emph{I can say a restaurant, a bakery, a cinema, a theatre and a town hall.}

Complete the last lesson of chapter 105: I can say a restaurant, a bakery, a cinema, a theatre and a town hall.
\end{chapteropening}

\section[das Restaurant]{das Restaurant — a restaurant}

\label{lesson:GE-C105-restaurant}



\begin{quote}
Before the new one: say the German for a hotel, then the German for the post office; the post.

An evening at the restaurant: say \emph{Guten Abend}, then \emph{der Abend}.
\end{quote}

\subsection*{You'll want to know: das Restaurant}
\textbf{das Restaurant} — \textquotedblleft{}a restaurant\textquotedblright{}.

\begin{cousinweb}
\textbf{Restaurant} $\leftarrow$ French \emph{restaurant}, \textquotedblleft{}restoring\textquotedblright{}, from \emph{restaurer}, \textquotedblleft{}to restore\textquotedblright{} $\leftarrow$ Latin \emph{restaurāre} (EtymDB 2.1) — borrowed: a place that restores you.
\end{cousinweb}

\begin{grammarlens}[title={What you've built}]
One more word for this chapter.
\end{grammarlens}

\subsection*{Your turn}
\begin{itemize}
\raggedright
  \item \emph{Say it:} \emph{das Restaurant}
  \item \emph{Say it:} \emph{das Restaurant}, once more
  \item \emph{Say it:} say \emph{das Restaurant}
  \item \emph{Recall:} say the German for a hotel, then the German for the post office; the post, then say \emph{das Restaurant} again
\end{itemize}

\begin{tcolorbox}[breakable,colback=teal!4,colframe=teal!35!black,title={Before you move on}]
What does das Restaurant mean? (\textquotedblleft{}A restaurant\textquotedblright{}.) Say it once more.
\end{tcolorbox}
\section[die Bäckerei]{die Bäckerei — a bakery}

\label{lesson:GE-C105-baeckerei}



\begin{quote}
Before the new one: say the German for the post office; the post, then the German for a restaurant.
\end{quote}

\subsection*{You'll want to know: die Bäckerei}
\textbf{die Bäckerei} — \textquotedblleft{}a bakery\textquotedblright{}.

\begin{cousinweb}
\textbf{Bäckerei} is the word for \textquotedblleft{}baker\textquotedblright{} + \emph{-ei}, the ending that names a trade's workplace (EtymDB 2.1). The baker is the one who bakes.
\end{cousinweb}

\begin{grammarlens}[title={What you've built}]
One more word for this chapter.
\end{grammarlens}

\subsection*{Your turn}
\begin{itemize}
\raggedright
  \item \emph{Say it:} \emph{die Bäckerei}
  \item \emph{Say it:} \emph{die Bäckerei}, once more
  \item \emph{Say it:} say \emph{die Bäckerei}
  \item \emph{Recall:} say the German for the post office; the post, then the German for a restaurant, then say \emph{die Bäckerei} again
\end{itemize}

\begin{tcolorbox}[breakable,colback=teal!4,colframe=teal!35!black,title={Before you move on}]
What does die Bäckerei mean? (\textquotedblleft{}A bakery\textquotedblright{}.) Say it once more.
\end{tcolorbox}
\section[das Kino]{das Kino — a cinema}

\label{lesson:GE-C105-kino}



\begin{quote}
Before the new one: say the German for a restaurant, then the German for a bakery.
\end{quote}

\subsection*{You'll want to know: das Kino}
\textbf{das Kino} — \textquotedblleft{}a cinema\textquotedblright{}.

\begin{cousinweb}
\textbf{Kino} is short for the word for the cinematograph, a modern coinage from Greek \emph{kínēma}, \textquotedblleft{}movement\textquotedblright{}, and \emph{gráphein}, \textquotedblleft{}to write\textquotedblright{}: the machine that records movement. EtymDB 2.1 has no source for \emph{Kino}.
\end{cousinweb}

\begin{grammarlens}[title={What you've built}]
One more word for this chapter.
\end{grammarlens}

\subsection*{Your turn}
\begin{itemize}
\raggedright
  \item \emph{Say it:} \emph{das Kino}
  \item \emph{Say it:} \emph{das Kino}, once more
  \item \emph{Say it:} say \emph{das Kino}
  \item \emph{Recall:} say the German for a restaurant, then the German for a bakery, then say \emph{das Kino} again
\end{itemize}

\begin{tcolorbox}[breakable,colback=teal!4,colframe=teal!35!black,title={Before you move on}]
What does das Kino mean? (\textquotedblleft{}A cinema\textquotedblright{}.) Say it once more.
\end{tcolorbox}
\section[das Theater]{das Theater — a theatre}

\label{lesson:GE-C105-theater}



\begin{quote}
Before the new one: say the German for a bakery, then the German for a cinema.
\end{quote}

\subsection*{You'll want to know: das Theater}
\textbf{das Theater} — \textquotedblleft{}a theatre\textquotedblright{}.

\begin{cousinweb}
\textbf{Theater} $\leftarrow$ Greek \emph{théatron}, \textquotedblleft{}a place for viewing\textquotedblright{}, from \emph{theáomai}, \textquotedblleft{}to watch\textquotedblright{} (EtymDB 2.1) — borrowed through Latin \emph{theātrum}.
\end{cousinweb}

\begin{grammarlens}[title={What you've built}]
One more word for this chapter.
\end{grammarlens}

\subsection*{Your turn}
\begin{itemize}
\raggedright
  \item \emph{Say it:} \emph{das Theater}
  \item \emph{Say it:} \emph{das Theater}, once more
  \item \emph{Say it:} say \emph{das Theater}
  \item \emph{Recall:} say the German for a bakery, then the German for a cinema, then say \emph{das Theater} again
\end{itemize}

\begin{tcolorbox}[breakable,colback=teal!4,colframe=teal!35!black,title={Before you move on}]
What does das Theater mean? (\textquotedblleft{}A theatre\textquotedblright{}.) Say it once more.
\end{tcolorbox}
\section[das Rathaus]{das Rathaus — a town hall}

\label{lesson:GE-C105-rathaus}



\begin{quote}
Before the new one: say the German for a cinema, then the German for a theatre.
\end{quote}

\subsection*{You'll want to know: das Rathaus}
\textbf{das Rathaus} — \textquotedblleft{}a town hall\textquotedblright{}.

\begin{cousinweb}
\textbf{Rathaus} is built from the word for \textquotedblleft{}council\textquotedblright{}, Old High German \emph{rāt} $\leftarrow$ Proto-Germanic \emph{*rēdaz}, \textquotedblleft{}advice\textquotedblright{} (EtymDB 2.1), and the word for \textquotedblleft{}house\textquotedblright{}. English \emph{rede}, \textquotedblleft{}counsel\textquotedblright{}, is a cousin.
\end{cousinweb}

\begin{grammarlens}[title={What you've built}]
One more word for this chapter.
\end{grammarlens}

\subsection*{Your turn}
\begin{itemize}
\raggedright
  \item \emph{Say it:} \emph{das Rathaus}
  \item \emph{Say it:} \emph{das Rathaus}, once more
  \item \emph{Say it:} say \emph{das Rathaus}
  \item \emph{Recall:} say the German for a cinema, then the German for a theatre, then say \emph{das Rathaus} again
\end{itemize}

\begin{tcolorbox}[breakable,colback=teal!4,colframe=teal!35!black,title={Before you move on}]
What does das Rathaus mean? (\textquotedblleft{}A town hall\textquotedblright{}.) Say it once more.
\end{tcolorbox}
